% Generated draft: requires mathematical and source review.
\documentclass[11pt]{article}
\usepackage[margin=25mm]{geometry}
\usepackage{fontspec}
\setmainfont{Latin Modern Roman}
\usepackage{amsmath,amssymb,mathtools,unicode-math}
\setmathfont{Latin Modern Math}
\usepackage{enumitem,needspace,xurl,hyperref}
\setlength{\parindent}{0pt}
\setlength{\parskip}{5pt}
\setlength{\emergencystretch}{3em}
\newcommand{\N}{\mathbb{N}}
\newcommand{\Z}{\mathbb{Z}}
\newcommand{\Q}{\mathbb{Q}}
\newcommand{\R}{\mathbb{R}}
\newcommand{\C}{\mathbb{C}}
\newcommand{\F}{\mathbb{F}}
\newcommand{\A}{\mathbb{A}}
\newcommand{\PP}{\mathbb{P}}
\newcommand{\E}{\mathbb{E}}
\newcommand{\eps}{\varepsilon}
\newcommand{\dd}{\,\mathrm{d}}
\begin{document}
\section*{UnsolvedMath ID: 30005714}
\textbf{Source problem code:} catalogue entry
\subsection{Definitions and mathematical statement}
A topological space \(X\) is contractible if there exist \(x_0\in X\) and a continuous map \(H:X\times[0,1]\to X\) such that \(H(x,0)=x\) and \(H(x,1)=x_0\) for every \(x\in X\). For a group \(G\), an action \(\alpha:G\times X\to X\) satisfies \(\alpha(e,x)=x\) and \(\alpha(gh,x)=\alpha(g,\alpha(h,x))\). Its global fixed-point set is
\[
X^G=\{x\in X:\alpha(g,x)=x\text{ for every }g\in G\}.
\]
\textbf{Casacuberta--Dicks conjecture, as reconstructed in the supplied record.} For every finite group \(G\), every contractible two-dimensional complex \(X\), with no compactness assumption, and every action of \(G\) on \(X\) in the intended category, one has
\[
X^G\ne\varnothing.
\]
Thus the conclusion requires a single point fixed simultaneously by all elements of \(G\).

\textbf{Unresolved scope.} The supplied record does not specify the category of complexes or the admissible actions. In particular, it does not determine whether complexes are CW complexes or simplicial complexes, or whether actions must be cellular, simplicial, or merely continuous. The action identities above do not resolve these additional requirements. It also does not clarify whether ``two-dimensional'' means dimension exactly \(2\) or dimension at most \(2\). Consequently, the displayed assertion records the supplied target with these qualifications unresolved; it does not assert that all these interpretations are equivalent. Removing compactness permits noncompact complexes and does not impose noncompactness as an additional hypothesis.
\subsection{Short English statement}
The supplied Casacuberta--Dicks conjecture asks whether every finite group action on a contractible two-dimensional complex has a common fixed point, without assuming compactness. The precise types of complexes and actions require clarification. The supplied record labels the noncompact case open; this draft does not independently verify that status.
\subsection{Sources}
\begin{itemize}
\item Supplied local UnsolvedMath registry record, permanent ID \(30005714\), title \emph{Noncompact Casacuberta-Dicks Conjecture}, source problem code \texttt{catalogue entry}. The conjecture is attributed by name to Casacuberta--Dicks; no individual proposer is supplied.
\item Supplied original report citation: \emph{Geometric, Algebraic and Topological Combinatorics} (2024), Oberwolfach Reports, DOI 10.4171/owr/2023/58, pp.\ 3249--3318. \url{https://doi.org/10.4171/owr/2023/58}. The retained original passage reports that the noncompact case remains open, without giving the complete conjecture or its hypotheses.
\item Supplied literature source: \emph{Oberwolfach Report 9/2026}. \url{https://ems.press/content/serial-article-files/53607}.
\item Supplied literature source labelled ``original Oberwolfach report'': \url{https://doi.org/10.4171/owr/2023/58}.
\item Literature assessment dated exactly \texttt{2026-08-22}: the supplied curation labels the problem open, reports recent positive results for the compact or finite case, and says that the cited current report explicitly leaves the noncompact case open. This is supplied dated curation, not an independently verified literature assessment; the status and cited claims require independent verification.
\item The registry labels the statement \texttt{corrected\_verified} and says it was reconstructed and checked against the official Oberwolfach report. These are supplied provenance claims, not verification performed for this draft. The record gives proposed year \(2023\), while its report citation gives \(2024\) and its DOI contains \(2023\); these supplied details are retained without reconciliation.
\end{itemize}
\end{document}
